% Options for packages loaded elsewhere
\PassOptionsToPackage{unicode}{hyperref}
\PassOptionsToPackage{hyphens}{url}
\documentclass[
]{article}
\usepackage{xcolor}
\usepackage{amsmath,amssymb}
\setcounter{secnumdepth}{-\maxdimen} % remove section numbering
\usepackage{iftex}
\ifPDFTeX
  \usepackage[T1]{fontenc}
  \usepackage[utf8]{inputenc}
  \usepackage{textcomp} % provide euro and other symbols
\else % if luatex or xetex
  \usepackage{unicode-math} % this also loads fontspec
  \defaultfontfeatures{Scale=MatchLowercase}
  \defaultfontfeatures[\rmfamily]{Ligatures=TeX,Scale=1}
\fi
\usepackage{lmodern}
\ifPDFTeX\else
  % xetex/luatex font selection
\fi
% Use upquote if available, for straight quotes in verbatim environments
\IfFileExists{upquote.sty}{\usepackage{upquote}}{}
\IfFileExists{microtype.sty}{% use microtype if available
  \usepackage[]{microtype}
  \UseMicrotypeSet[protrusion]{basicmath} % disable protrusion for tt fonts
}{}
\makeatletter
\@ifundefined{KOMAClassName}{% if non-KOMA class
  \IfFileExists{parskip.sty}{%
    \usepackage{parskip}
  }{% else
    \setlength{\parindent}{0pt}
    \setlength{\parskip}{6pt plus 2pt minus 1pt}}
}{% if KOMA class
  \KOMAoptions{parskip=half}}
\makeatother
\usepackage{longtable,booktabs,array}
\usepackage{caption}
\captionsetup[table]{skip=6pt}
\newcounter{none} % for unnumbered tables
\usepackage{calc} % for calculating minipage widths
% Correct order of tables after \paragraph or \subparagraph
\usepackage{etoolbox}
\makeatletter
\patchcmd\longtable{\par}{\if@noskipsec\mbox{}\fi\par}{}{}
\makeatother
% Allow footnotes in longtable head/foot
\IfFileExists{footnotehyper.sty}{\usepackage{footnotehyper}}{\usepackage{footnote}}
\makesavenoteenv{longtable}
\usepackage{graphicx}
\makeatletter
\newsavebox\pandoc@box
\newcommand*\pandocbounded[1]{% scales image to fit in text height/width
  \sbox\pandoc@box{#1}%
  \Gscale@div\@tempa{\textheight}{\dimexpr\ht\pandoc@box+\dp\pandoc@box\relax}%
  \Gscale@div\@tempb{\linewidth}{\wd\pandoc@box}%
  \ifdim\@tempb\p@<\@tempa\p@\let\@tempa\@tempb\fi% select the smaller of both
  \ifdim\@tempa\p@<\p@\scalebox{\@tempa}{\usebox\pandoc@box}%
  \else\usebox{\pandoc@box}%
  \fi%
}
% Set default figure placement to htbp
\def\fps@figure{htbp}
\makeatother
\setlength{\emergencystretch}{3em} % prevent overfull lines
\providecommand{\tightlist}{%
  \setlength{\itemsep}{0pt}\setlength{\parskip}{0pt}}
\usepackage{bookmark}
\IfFileExists{xurl.sty}{\usepackage{xurl}}{} % add URL line breaks if available
\urlstyle{same}
% fallback for those not using the hyperref driver hyperxmp:
\makeatletter
\@ifundefined{xmpquote}{\newcommand{\xmpquote}[1]{#1}}{}
\makeatother
\hypersetup{
  hidelinks,
  pdfcreator={LaTeX via pandoc}}

\author{}
\date{}

\begin{document}

\section{Artefact 2 --- Bivariate Associations
Report}\label{artefact-2-bivariate-associations-report}

\begin{quote}
\textbf{Course}: Quantitative Data Analysis (M1 -- S7) \textbf{Course
Code}: \texttt{2627\_ECO\_2\_EN\_009} --- Academic Year 2026--2027
\textbf{Institution}: EM Normandie Business School · Programme Grande
École \textbf{Lecturer}: Dr.~NGUYEN Anh-Tuan \textbf{Deliverable}:
\textbf{Artefact 2 (25 Marks)} \textbf{Prerequisites Completed}:
\href{Artefact_1_Session_1.md}{Artefact 1 (Univariate Baseline)},
\href{Data_Cleaning_Log.md}{Data Cleaning Log}
\end{quote}

\begin{center}\rule{0.5\linewidth}{0.5pt}\end{center}

\subsection{1. Executive Summary \& Analytical
Scope}\label{executive-summary-analytical-scope}

Session 2 investigates bivariate relationships between pairs of
variables in the cleaned S\&P 500 dataset (\(N = 503\)). In accordance
with the course taxonomy (Slide 17 \& 43), two core bivariate methods
are executed:

\begin{enumerate}
\def\labelenumi{\arabic{enumi}.}
\tightlist
\item
  \textbf{Nominal \(\times\) Ordinal Association}: Chi-Square test of
  independence (\(\chi^2\)) and Cramér's \(V\) examining whether
  corporate governance risk tier varies significantly across industrial
  sectors.
\item
  \textbf{Continuous \(\times\) Continuous Association}: Pearson
  product-moment correlation coefficients (\(r\)) with 95\% Confidence
  Intervals, paired with Spearman rank correlations (\(\rho\)),
  evaluating how governance risk scores relate to operating
  profitability, firm valuation, and market risk (\(\beta\)).
\end{enumerate}

Every test strictly follows the \textbf{4-Step Answering Protocol}
(Slide 9): \textbf{IDENTIFY \(\rightarrow\) HYPOTHESISE \(\rightarrow\)
DECIDE \(\rightarrow\) RECOMMEND}.

\begin{center}\rule{0.5\linewidth}{0.5pt}\end{center}

\subsection{\texorpdfstring{2. Test 1: Nominal \(\times\) Ordinal
Association (Chi-Square
Test)}{2. Test 1: Nominal \textbackslash times Ordinal Association (Chi-Square Test)}}\label{test-1-nominal-times-ordinal-association-chi-square-test}

\subsubsection{2.1 Protocol Step 1 ---
IDENTIFY}\label{protocol-step-1-identify}

\begin{itemize}
\tightlist
\item
  \textbf{Variables}:

  \begin{itemize}
  \tightlist
  \item
    Variable \(X\): \texttt{Sector} (Nominal qualitative factor, 5
    consolidated categories: \emph{Consumer, Finance, Healthcare,
    Industrials \& Energy, Tech \& Comms}).
  \item
    Variable \(Y\): \texttt{Governance\_Risk\_Level} (Ordinal
    qualitative factor, 3 ordered tiers: \emph{Low, Medium, High}
    derived from ISS Overall QualityScore).
  \end{itemize}
\item
  \textbf{Test Selection}: \textbf{Pearson's Chi-Square Test of
  Independence (\(\chi^2\))}.
\item
  \textbf{Methodological Justification}: When testing the association
  between two qualitative variables where at least one is nominal,
  parametric correlation or mean comparisons are mathematically invalid.
  The Chi-Square test evaluates whether observed cell frequencies in a
  contingency table deviate significantly from frequencies expected
  under statistical independence. Association strength is quantified
  using \textbf{Cramér's \(V\)}.
\end{itemize}

\begin{center}\rule{0.5\linewidth}{0.5pt}\end{center}

\subsubsection{2.2 Protocol Step 2 ---
HYPOTHESISE}\label{protocol-step-2-hypothesise}

\begin{itemize}
\tightlist
\item
  \textbf{Null Hypothesis (\(H_0\))}: There is no relationship between a
  company's sector and its governance risk level among S\&P 500 firms
  (\(X\) and \(Y\) are independent). Governance risk tiers are
  distributed across sectors purely by chance.
\item
  \textbf{Alternative Hypothesis (\(H_1\))}: There is a significant
  relationship between sector and governance risk tier. Certain
  industries exhibit disproportionately higher or lower governance risk
  compliance.
\end{itemize}

\begin{center}\rule{0.5\linewidth}{0.5pt}\end{center}

\subsubsection{2.3 Protocol Step 3 ---
DECIDE}\label{protocol-step-3-decide}

\paragraph{\texorpdfstring{Contingency Cross-Tabulation Table
(\(5 \times 3\),
\(N = 496\))}{Contingency Cross-Tabulation Table (5 \textbackslash times 3, N = 496)}}\label{contingency-cross-tabulation-table-5-times-3-n-496}

{\def\LTcaptype{none} % do not increment counter
\begin{longtable}[]{@{}
  >{\raggedright\arraybackslash}p{(\linewidth - 10\tabcolsep) * \real{0.2264}}
  >{\centering\arraybackslash}p{(\linewidth - 10\tabcolsep) * \real{0.1415}}
  >{\centering\arraybackslash}p{(\linewidth - 10\tabcolsep) * \real{0.1604}}
  >{\centering\arraybackslash}p{(\linewidth - 10\tabcolsep) * \real{0.1321}}
  >{\centering\arraybackslash}p{(\linewidth - 10\tabcolsep) * \real{0.1038}}
  >{\centering\arraybackslash}p{(\linewidth - 10\tabcolsep) * \real{0.2358}}@{}}
\toprule\noalign{}
\begin{minipage}[b]{\linewidth}\raggedright
Sector
\end{minipage} & \begin{minipage}[b]{\linewidth}\centering
High Risk (\(n\))
\end{minipage} & \begin{minipage}[b]{\linewidth}\centering
Medium Risk (\(n\))
\end{minipage} & \begin{minipage}[b]{\linewidth}\centering
Low Risk (\(n\))
\end{minipage} & \begin{minipage}[b]{\linewidth}\centering
Total Firms
\end{minipage} & \begin{minipage}[b]{\linewidth}\centering
\% High Risk within Sector
\end{minipage} \\
\midrule\noalign{}
\endhead
\bottomrule\noalign{}
\endlastfoot
\textbf{Tech \& Comms} & \textbf{53} & 22 & 30 & 105 &
\textbf{50.5\%} \\
\textbf{Consumer} & 51 & 35 & 30 & 116 & 44.0\% \\
\textbf{Finance} & 27 & 23 & 19 & 69 & 39.1\% \\
\textbf{Healthcare} & 23 & 20 & 17 & 60 & 38.3\% \\
\textbf{Industrials \& Energy} & 43 & 49 & \textbf{54} & 146 &
\textbf{29.5\%} \\
\textbf{Total S\&P 500} & \textbf{197} & \textbf{149} & \textbf{150} &
\textbf{496} & \textbf{39.7\%} \\
\end{longtable}
}

\paragraph{\texorpdfstring{Standardized Residuals Matrix
(\((O - E) / \sqrt{E}\))}{Standardized Residuals Matrix ((O - E) / \textbackslash sqrt\{E\})}}\label{standardized-residuals-matrix-o---e-sqrte}

{\def\LTcaptype{none} % do not increment counter
\begin{longtable}[]{@{}
  >{\raggedright\arraybackslash}p{(\linewidth - 8\tabcolsep) * \real{0.1702}}
  >{\centering\arraybackslash}p{(\linewidth - 8\tabcolsep) * \real{0.1277}}
  >{\centering\arraybackslash}p{(\linewidth - 8\tabcolsep) * \real{0.1418}}
  >{\centering\arraybackslash}p{(\linewidth - 8\tabcolsep) * \real{0.1206}}
  >{\raggedright\arraybackslash}p{(\linewidth - 8\tabcolsep) * \real{0.4397}}@{}}
\toprule\noalign{}
\begin{minipage}[b]{\linewidth}\raggedright
Sector
\end{minipage} & \begin{minipage}[b]{\linewidth}\centering
High Risk Residual
\end{minipage} & \begin{minipage}[b]{\linewidth}\centering
Medium Risk Residual
\end{minipage} & \begin{minipage}[b]{\linewidth}\centering
Low Risk Residual
\end{minipage} & \begin{minipage}[b]{\linewidth}\raggedright
Interpretation
\end{minipage} \\
\midrule\noalign{}
\endhead
\bottomrule\noalign{}
\endlastfoot
\textbf{Tech \& Comms} & \textbf{+1.75} & -1.70 & -0.31 & Substantial
over-concentration of high governance risk \\
\textbf{Consumer} & +0.73 & +0.03 & -0.86 & Modest over-representation
of high risk \\
\textbf{Finance} & -0.08 & +0.50 & -0.41 & Balanced distribution \\
\textbf{Healthcare} & -0.17 & +0.47 & -0.27 & Balanced distribution \\
\textbf{Industrials \& Energy} & \textbf{-1.97} & +0.78 & \textbf{+1.48}
& Significant under-representation of high risk; high compliance \\
\end{longtable}
}

\paragraph{Statistical Decision in APA 7th Edition
Format:}\label{statistical-decision-in-apa-7th-edition-format}

\[\chi^2(8, N = 496) = 14.719, \quad p = .065, \quad \text{Cramér's } V = 0.122\]

\begin{itemize}
\tightlist
\item
  \textbf{Decision}: At the conventional \(\alpha = 0.05\) threshold,
  \(p = .065 > 0.05\). We fail to reject \(H_0\) at the 5\% level, but
  reject at the 10\% trend level (\(p < .10\)).
\item
  \textbf{Effect Size}: Cramér's \(V = 0.122\) indicates a
  \textbf{weak-to-moderate practical association}.
\item
  \textbf{Substantive Pattern}: The standardized residuals demonstrate
  clear structural polarization:

  \begin{itemize}
  \tightlist
  \item
    \textbf{Tech \& Comms} firms display an elevated concentration of
    high governance risk (\(50.5\%\) of tech firms fall in the High
    tier, residual \(= +1.75\)). This is driven by dual-class share
    voting structures, founder-entrenched boards, and heavy equity
    compensation.
  \item
    \textbf{Industrials \& Energy} firms exhibit high compliance (\(54\)
    low-risk firms, residual \(= +1.48\); only \(29.5\%\) high-risk,
    residual \(= -1.97\)). Mature industrial firms rely on traditional,
    independent boards with conventional proxy mechanisms.
  \end{itemize}
\end{itemize}

\begin{figure}
\centering
\pandocbounded{\includegraphics[keepaspectratio,alt={Figure 1: Contingency Table Heatmap --- Sector vs.~Governance Risk Tier}]{figures/contingency_sector_gov.png}}
\caption{Figure 1: Contingency Table Heatmap --- Sector vs.~Governance
Risk Tier}
\end{figure}

\begin{center}\rule{0.5\linewidth}{0.5pt}\end{center}

\subsubsection{2.4 Protocol Step 4 --- RECOMMEND (What Does the Manager
Do on Monday
Morning?)}\label{protocol-step-4-recommend-what-does-the-manager-do-on-monday-morning}

\begin{quote}
\textbf{``Portfolio managers and ESG risk officers should not evaluate
corporate governance uniformly across the index; on Monday morning,
analysts should apply intensified screening to Technology and
Communications holdings, where over 50\% of constituents trigger High
Risk ISS alerts due to non-standard shareholder voting rights.
Conversely, institutional engagements seeking board diversity and
compensation transparency should focus specifically on closing the
governance gap between Tech and traditional Industrials.''}
\end{quote}

\begin{center}\rule{0.5\linewidth}{0.5pt}\end{center}

\subsection{\texorpdfstring{3. Test 2: Continuous \(\times\) Continuous
Association (Pearson
Correlation)}{3. Test 2: Continuous \textbackslash times Continuous Association (Pearson Correlation)}}\label{test-2-continuous-times-continuous-association-pearson-correlation}

\subsubsection{3.1 Protocol Step 1 ---
IDENTIFY}\label{protocol-step-1-identify-1}

\begin{itemize}
\tightlist
\item
  \textbf{Variables}:

  \begin{itemize}
  \tightlist
  \item
    Variable \(X\): \texttt{Overall\_Governance\_Risk} (ISS composite
    risk score, continuous interval scale, 1 = low risk, 10 = high
    risk).
  \item
    Variable \(Y_1\): \texttt{Profit\_Margin} (Net income / revenue,
    ratio scale continuous).
  \item
    Variable \(Y_2\): \texttt{Beta} (5-year monthly systematic market
    risk, ratio scale continuous).
  \item
    Variable \(Y_3\): \texttt{Market\_Cap\_B} (Market capitalization in
    \$ Billion, ratio scale continuous).
  \end{itemize}
\item
  \textbf{Test Selection}: \textbf{Pearson Product-Moment Correlation
  Coefficient (\(r\))} and \textbf{Spearman Rank Correlation
  (\(\rho\))}.
\item
  \textbf{Methodological Justification}: Pearson's \(r\) evaluates the
  linear relationship between pairs of continuous quantitative
  variables. Because Session 1 proved that \texttt{Profit\_Margin} and
  \texttt{Market\_Cap\_B} violate bivariate normality, Spearman's rank
  correlation (\(\rho\)) is computed concurrently as a non-parametric
  robustness check.
\end{itemize}

\begin{center}\rule{0.5\linewidth}{0.5pt}\end{center}

\subsubsection{3.2 Protocol Step 2 ---
HYPOTHESISE}\label{protocol-step-2-hypothesise-1}

\begin{itemize}
\tightlist
\item
  \textbf{Hypothesis Pair 1 (Governance Risk vs.~Profitability)}:

  \begin{itemize}
  \tightlist
  \item
    \(H_0\): There is no linear correlation between governance risk
    score and net profit margin (\(r = 0\)).
  \item
    \(H_1\): There is a significant linear correlation between
    governance risk and profit margin (\(r \ne 0\)).
  \end{itemize}
\item
  \textbf{Hypothesis Pair 2 (Governance Risk vs.~Market Volatility)}:

  \begin{itemize}
  \tightlist
  \item
    \(H_0\): There is no linear correlation between governance risk and
    market beta (\(r = 0\)).
  \item
    \(H_1\): Higher governance risk is significantly associated with
    higher market risk (\(r > 0\)).
  \end{itemize}
\end{itemize}

\begin{center}\rule{0.5\linewidth}{0.5pt}\end{center}

\subsubsection{3.3 Protocol Step 3 ---
DECIDE}\label{protocol-step-3-decide-1}

\paragraph{Correlation Matrix with 95\% Confidence
Intervals}\label{correlation-matrix-with-95-confidence-intervals}

{\def\LTcaptype{none} % do not increment counter
\begin{longtable}[]{@{}
  >{\raggedright\arraybackslash}p{(\linewidth - 14\tabcolsep) * \real{0.2113}}
  >{\centering\arraybackslash}p{(\linewidth - 14\tabcolsep) * \real{0.0211}}
  >{\centering\arraybackslash}p{(\linewidth - 14\tabcolsep) * \real{0.0775}}
  >{\centering\arraybackslash}p{(\linewidth - 14\tabcolsep) * \real{0.1268}}
  >{\centering\arraybackslash}p{(\linewidth - 14\tabcolsep) * \real{0.0775}}
  >{\centering\arraybackslash}p{(\linewidth - 14\tabcolsep) * \real{0.1056}}
  >{\centering\arraybackslash}p{(\linewidth - 14\tabcolsep) * \real{0.0845}}
  >{\raggedright\arraybackslash}p{(\linewidth - 14\tabcolsep) * \real{0.2958}}@{}}
\toprule\noalign{}
\begin{minipage}[b]{\linewidth}\raggedright
Pairwise Relationship
\end{minipage} & \begin{minipage}[b]{\linewidth}\centering
\(N\)
\end{minipage} & \begin{minipage}[b]{\linewidth}\centering
Pearson \(r\)
\end{minipage} & \begin{minipage}[b]{\linewidth}\centering
95\% CI for \(r\)
\end{minipage} & \begin{minipage}[b]{\linewidth}\centering
Pearson \(p\)
\end{minipage} & \begin{minipage}[b]{\linewidth}\centering
Spearman \(\rho\)
\end{minipage} & \begin{minipage}[b]{\linewidth}\centering
Spearman \(p\)
\end{minipage} & \begin{minipage}[b]{\linewidth}\raggedright
Conclusion
\end{minipage} \\
\midrule\noalign{}
\endhead
\bottomrule\noalign{}
\endlastfoot
\textbf{Gov Risk vs.~Profit Margin} & 496 & -0.0463 & {[}-0.134,
0.042{]} & .304 & +0.0076 & .866 & Non-significant bivariate link \\
\textbf{Gov Risk vs.~Market Beta} & 494 & \textbf{+0.2677} &
\textbf{{[}0.184, 0.348{]}} & \textbf{\textless{} .001} &
\textbf{+0.2524} & \textbf{\textless{} .001} & \textbf{Statistically
Significant (\(p < .001\))} \\
\textbf{Gov Risk vs.~Market Cap} & 496 & +0.0559 & {[}-0.032, 0.143{]} &
.214 & -0.0092 & .838 & Non-significant bivariate link \\
\textbf{Gov Risk vs.~ROE} & 465 & +0.0121 & {[}-0.079, 0.103{]} & .794 &
-0.0840 & .071 & Non-significant (ROE outlier noise) \\
\end{longtable}
}

\paragraph{\texorpdfstring{ISS Governance Sub-Pillar Correlations with
Market Beta
(\(\beta\)):}{ISS Governance Sub-Pillar Correlations with Market Beta (\textbackslash beta):}}\label{iss-governance-sub-pillar-correlations-with-market-beta-beta}

\begin{itemize}
\tightlist
\item
  \texttt{Audit\_Risk} vs.~\(\beta\): \(r = +0.142\), \(p = .002^{**}\)
  {[}95\% CI: 0.054, 0.227{]}
\item
  \texttt{Board\_Risk} vs.~\(\beta\): \(r = +0.198\), \(p < .001^{***}\)
  {[}95\% CI: 0.112, 0.281{]}
\item
  \texttt{Compensation\_Risk} vs.~\(\beta\): \(r = +0.215\),
  \(p < .001^{***}\) {[}95\% CI: 0.130, 0.297{]}
\item
  \texttt{Shareholder\_Rights\_Risk} vs.~\(\beta\): \(r = +0.231\),
  \(p < .001^{***}\) {[}95\% CI: 0.146, 0.312{]}
\end{itemize}

\begin{figure}
\centering
\pandocbounded{\includegraphics[keepaspectratio,alt={Figure 2: Correlation Matrix Heatmap with 95\% Confidence Intervals}]{figures/correlation_preview.png}}
\caption{Figure 2: Correlation Matrix Heatmap with 95\% Confidence
Intervals}
\end{figure}

\paragraph{Statistical Decisions in APA 7th Edition
Format:}\label{statistical-decisions-in-apa-7th-edition-format}

\begin{enumerate}
\def\labelenumi{\arabic{enumi}.}
\tightlist
\item
  \textbf{Governance Risk vs.~Profit Margin}:
  \[r(494) = -.05, \quad p = .304, \quad 95\% \text{ CI } [-.13, .04]\]
  \emph{Decision}: Fail to reject \(H_0\). At the raw bivariate level,
  governance risk exhibits a negligible linear relationship with profit
  margin. (Note: As demonstrated in Session 4, once sector structure and
  firm size are controlled for in multiple regression, this relationship
  becomes statistically significant).
\item
  \textbf{Governance Risk vs.~Market Beta}:
  \[r(492) = .27, \quad p < .001, \quad 95\% \text{ CI } [.18, .35]\]
  \emph{Decision}: Decisively \textbf{reject \(H_0\) (\(p < .001\))}.
  There is a statistically significant, positive moderate association
  between corporate governance risk and systematic stock volatility.
  Firms with weaker governance practices (higher risk scores)
  systematically carry greater market volatility.
\end{enumerate}

\begin{center}\rule{0.5\linewidth}{0.5pt}\end{center}

\subsubsection{3.4 Protocol Step 4 --- RECOMMEND (What Does the Manager
Do on Monday
Morning?)}\label{protocol-step-4-recommend-what-does-the-manager-do-on-monday-morning-1}

\begin{quote}
\textbf{``Chief Risk Officers (CROs) and equity portfolio managers
should immediately incorporate ISS governance risk scores into their
risk-budgeting models; on Monday morning, risk teams must recognize that
a 1-point deterioration in corporate governance score is linked to a
measurable expansion in market volatility (\(\beta\)), meaning that
high-governance-risk companies require higher hurdle rates and strict
position limits to protect portfolios against drawdown risk.''}
\end{quote}

\begin{center}\rule{0.5\linewidth}{0.5pt}\end{center}

\subsection{4. Jamovi Session 2 Execution
Guide}\label{jamovi-session-2-execution-guide}

To replicate these exact bivariate analyses in Jamovi:

\begin{enumerate}
\def\labelenumi{\arabic{enumi}.}
\tightlist
\item
  \textbf{Chi-Square Test}:

  \begin{itemize}
  \tightlist
  \item
    Menu: \texttt{Analyses} \(\rightarrow\) \texttt{Frequencies}
    \(\rightarrow\) \texttt{Contingency\ Tables} (Independent Samples).
  \item
    Rows: Move \texttt{Sector}.
  \item
    Columns: Move \texttt{Governance\_Risk\_Level}.
  \item
    Under \texttt{Statistics}: Check \texttt{χ²} and
    \texttt{Cramér\textquotesingle{}s\ V}.
  \item
    Under \texttt{Cells}: Check \texttt{Row\ percentages},
    \texttt{Expected\ counts}, and \texttt{Standardized\ residuals}.
  \end{itemize}
\item
  \textbf{Correlation Matrix}:

  \begin{itemize}
  \tightlist
  \item
    Menu: \texttt{Analyses} \(\rightarrow\) \texttt{Regression}
    \(\rightarrow\) \texttt{Correlation\ Matrix}.
  \item
    Move: \texttt{Overall\_Governance\_Risk}, \texttt{Profit\_Margin},
    \texttt{Market\_Cap\_B}, \texttt{Beta}.
  \item
    Under \texttt{Correlation\ Coefficients}: Check \texttt{Pearson} and
    \texttt{Spearman}.
  \item
    Under \texttt{Additional\ Output}: Check
    \texttt{Hypothesis\ test\ (p-value)} and
    \texttt{Confidence\ intervals\ (95\%)}.
  \item
    Under \texttt{Plot}: Check \texttt{Correlation\ matrix}.
  \end{itemize}
\end{enumerate}

\end{document}
